\section{Mary's Cooperation: Co-redemptrix, Mediatrix, and Advocate}

\dogmathesis{Mary's cooperation in salvation is real because grace perfects created freedom and because God willed the Incarnation to include her maternal consent, discipleship, and charity. It is wholly received because Christ alone accomplishes redemption as the incarnate Son and one perfect mediator. The current Magisterium therefore commends precise language of maternal cooperation and intercession, judges ``Co-redemptrix'' always inappropriate, limits ``Mediatrix'' to a distant and subordinate analogy, and receives ``Advocate'' only as maternal intercession within---never against---Christ's mediation.}

This chapter begins with doctrine rather than vocabulary. A title cannot rescue a defective account, and rejecting an ambiguous title need not deny the reality its advocates intended to protect. The governing sources are 1 Timothy 2:5--6; Vatican II, LG 56, 58, and 60--62; CCC 967--970; and the Dicastery for the Doctrine of the Faith's doctrinal note \emph{Mater Populi Fidelis} (4 November 2025), approved by Pope Leo XIV. MPF explicitly distinguishes legitimate Marian cooperation from formulas that compromise the unique mediation of Christ.

\subsection{The one Redeemer and perfect mediator}

There is one God and one mediator between God and humanity, the man Christ Jesus, who gave himself as ransom for all (1 Tim 2:5--6). The uniqueness is not merely that Christ is first among several comparable mediators. He is the eternal Son made man, the one subject who belongs fully to the divine and human orders and reconciles in his own person. His obedient life and once-for-all self-offering accomplish redemption; his risen priesthood communicates its fruits (Heb 7:23--28; 9:11--15; 10:10--14).

St.~Thomas distinguishes this perfect mediation from subordinate participation. A mediator joins extremes by acting between them. Christ alone unites God and humanity perfectly, reconciling through satisfaction and communicating divine gifts. Others may be called mediators \emph{secundum quid} insofar as they dispose or minister, but only dependently (ST III, q.26, aa.1--2). The distinction does not reserve all meaningful agency to Christ and make creatures inert. It identifies an unshareable source and a shareable participation.

Vatican II applies the same logic: Mary's maternal role ``in no way obscures or diminishes'' Christ's unique mediation but shows its power; no creature can be counted with the incarnate Word and Redeemer, while Christ's one goodness is really shared among creatures in diverse ways (LG 60, 62). Participation is possible because Christ is sufficient, not because he lacks something.

\subsection{Grace-enabled cooperation}

Mary cooperates first by receiving. At the Annunciation the Father gives vocation, the Spirit overshadows, and the Son assumes; Mary believes and consents (Luke 1:26--38). Vatican II says the Father willed that the predestined mother's consent should precede the Incarnation, and that she freely cooperated through faith and obedience in human salvation (LG 56). ``Freely'' and ``by grace'' are not rivals. God causes the created act according to the nature of a free cause; he does not bargain with an independent sovereignty.

Her cooperation unfolds through maternity and discipleship: visitation, birth, nurture, presentation, flight, hidden life, Cana, fidelity to the word, and presence at the Cross. At Calvary she does not cause the hypostatic union, consecrate the victim, exercise Christ's priesthood, add a second sacrifice, or supply merits missing from the Cross. She consents in faith to the immolation of the victim born of her and suffers maternally with him (LG 58). Her compassion is a grace-enabled human participation in the event whose redemptive efficacy belongs to Christ.

It is useful, with care, to distinguish the objective accomplishment of redemption from the reception and application of its fruits. Christ alone merits and effects reconciliation in the strict sense. The Church, sacraments, preaching, prayer, saints, and every Christian witness participate instrumentally or dispositively in communicating what he accomplished. Mary has a unique maternal place in both the historical advent of the Redeemer and the Church's reception of grace, but she never becomes a second objective source.

\subsection{The title ``Co-redemptrix''}

MPF 17--21 traces the title's emergence and the shift from earlier language of Mary as \emph{Redemptrix} to \emph{Co-redemptrix}. The Latin prefix \emph{co-} can mean ``with'' rather than ``equal,'' and responsible proponents often intend subordinate cooperation. Some popes used the expression or cognate formulations in non-defining addresses, and many saints and theologians used it to honor Mary's fiat and compassion. That history should be reported, neither suppressed nor inflated.

The title nevertheless carries an objective linguistic and doctrinal risk. In ordinary hearing ``co-redeemer'' suggests a second redeemer alongside Christ or a division of the saving work. Repeated explanation becomes necessary precisely where a title ought to illuminate. The preparatory work of Vatican II considered Marian cooperation but deliberately did not use or define \emph{Co-redemptrix}; the Council instead placed Mary's role under Christ's unique mediation. Cardinal Joseph Ratzinger later judged the formula too remote from Scripture and the Fathers and too liable to misunderstanding; Pope Francis repeatedly rejected it.

The DDF's present conclusion is categorical: ``it is always inappropriate to use the title Co-redemptrix'' (MPF 22). The reason is Christological, not a denial of Mary. The title risks obscuring the unique saving mediation of Christ and creates confusion where precise exposition is required. A Catholic reference written after this judgment should use the term only to document its history or explain the ruling, not recommend it, propose it as a fifth dogma, or suggest that petition campaigns have accumulated definitional force.

The preferred language is the Council's: Mary freely cooperated, served the mystery of redemption, was united with maternal heart to her Son, and participates maternally in the communication of grace. These phrases state more rather than less because they identify the acts and preserve their dependence.

\subsection{The title ``Mediatrix''}

``Mediatrix'' has a different history and judgment. MPF 23--26 traces the concept in Eastern Christian usage beginning in the sixth century, its more frequent Western use from the twelfth century, and its formulation as a doctrinal thesis in the seventeenth. In modern times Cardinal D\'esir\'e-Joseph Mercier promoted a definition of universal Marian mediation. Benedict XV granted a proper Mass and office for Mary Mediatrix of All Graces to Belgium in 1921, with later grants to other dioceses and religious orders; a limited liturgical concession was not a universal dogmatic definition. The theological commission and Council Fathers at Vatican II chose not to issue the proposed new definition.

LG 62 nevertheless receives \emph{Mediatrix} among traditional invocations, immediately adding that these titles are to be understood so that they neither take away from nor add anything to the dignity and efficacy of Christ the one Mediator. The title can therefore be orthodox in a carefully delimited analogical sense. Mary intercedes, maternally accompanies, exemplifies faith, and disposes believers to receive the gifts of her Son. Her mediation is participated, subordinate, and dependent---not an intermediate layer through which God is metaphysically unable to reach creatures.

MPF 27--29 and 53--55 insist that the analogy is distant. Christ's mediation is unique in being constitutive and perfect. Mary belongs on the redeemed human side; everything in her is received from his fullness. She can help persons encounter Christ, but she does not stand between the soul and Christ as a barrier, price, or necessary independent conduit. The most authentic Marian function repeats Cana: noticing need, interceding, and directing servants to obey Jesus.

\subsection{``Mediatrix of All Graces'' and the causality of grace}

The expanded title ``Mediatrix of All Graces'' is especially problematic if ``all'' implies that every divine gift must pass through Mary's causal possession or conscious distribution. MPF 45--55 and 67--70 reject an account in which Mary has a separate treasury or becomes a universal dispenser between God and each soul. The Trinity alone is the source of uncreated and created grace. Sanctifying grace is a participated divine life immediately caused by God in the soul; no creature, not even the glorified Mother of God, inserts itself into that innermost divine communication.

Thomas's sacramental theology gives the metaphysical rule: God alone enters the soul and causes grace as principal agent; ministers act instrumentally under Christ (ST III, q.64, a.1). His theology of merit likewise reserves strict merit of another's first grace to Christ, while allowing prayer and charity a subordinate fittingness under God's gracious ordination (ST I--II, q.114, a.6). Mary's intercession can therefore obtain actual helps according to providence; she can prepare, encourage, exemplify, and pray. She does not manufacture sanctifying grace, own it apart from God, or decide its distribution by a parallel sovereignty.

The rejection of a necessary chain is also Trinitarian. The Father does not delegate mercy because he is distant; the Son is not reluctant until his mother persuades him; the Spirit does not need a created bridge to indwell. In the economy God freely gives creatures roles, and their intercession is itself a fruit of his one generous will. The glorified Mary's prayer is powerful because God perfects communion, not because divine action is weak without her.

Consequently, a sentence such as ``every grace comes through Mary'' requires so much qualification that it is unsafe as an unqualified doctrinal axiom. At most, it can signify that every grace comes through Christ's Incarnation, to which Mary freely consented, or that her maternal intercession has a universal ecclesial scope by divine gift. It cannot mean efficient causality over each grace or a revealed universal rule that God never acts without a distinct Marian transaction.

\subsection{The title ``Advocate''}

``Advocate'' is constructive when it names maternal intercession within the communion of saints. Scripture gives the primary title to Christ, our advocate with the Father (1 John 2:1), and calls the Holy Spirit the Paraclete or Advocate (John 14:16, 26; 15:26; 16:7). Applied to Mary, the term is analogical and creaturely. She prays; she does not become another divine Paraclete or legal defense against Christ.

Cana supplies the controlling evangelical icon (John 2:1--11). Mary perceives the lack of wine, states the need without dictating the sign, trusts Jesus, and orders the servants to him. The sign comes from Christ and manifests his glory. In heaven the saints remain members of one body and their prayer is intensified, not abolished (Rev 5:8; ST II--II, q.83, a.11; CCC 956). Mary's intercession has a singular maternal character because of her relation to Christ and the Church.

LG 62 invokes her as Advocate, Auxiliatrix, Adjutrix, and Mediatrix under the explicit Christological safeguard. MPF 36--44 positively expounds maternal cooperation and intercession, and MPF 54 can speak of Mary as an ``advocate of grace'' when the phrase means prayer and maternal assistance. St.~Irenaeus's description of obedient Mary as advocate of Eve (\emph{Against Heresies} V.19.1) supplies an ancient recapitulative root: the New Eve's faith answers the old Eve's unbelief through the one new Adam.

The title becomes corrupt when devotional rhetoric portrays the Father as severe, Christ as judge, and Mary as the merciful party who restrains their wrath. Father, Son, and Spirit possess one divine will and one mercy. Christ himself intercedes and has given his life for sinners (Rom 8:31--34). Mary does not save anyone \emph{from} Jesus. Her advocacy is an expression of his Body's communion and his gift.

\subsection{Intercession and divine omnipotence}

God does not need creaturely intercession because of ignorance, incapacity, or emotional reluctance. He wills it because omnipotence can give creatures real participation without surrendering primacy. When Christians pray for one another, their prayer does not supply information or overpower Providence; it is one of the secondary causes eternally included in Providence. The saints' heavenly intercession has the same dependent structure.

This prevents two reductions. A mechanistic reduction imagines that if God already wills the good, prayer is pointless. A competitive reduction imagines prayer works only by changing a divine refusal. Classical providence includes the freely offered petition and the granted gift within one wise order. Mary's advocacy is neither decorative nor coercive: it is a real created act whose existence and efficacy come from God.

\subsection{A doctrinal test for Marian language}

Before using a title or formula, ask:

\begin{enumerate}
  \item Does it confess the Father's initiative, the Son's sufficient redemption, and the Spirit's immediate sanctifying action?
  \item Does it keep Mary on the created and redeemed side of the Creator--creature distinction?
  \item Does it identify her actual mode of cooperation---faith, consent, maternity, discipleship, compassion, example, or intercession---rather than invoke an unexplained power?
  \item Does it distinguish Christ's perfect mediation from distant, participated mediation?
  \item Could ordinary hearers reasonably understand the words without receiving a false two-redeemer or universal-dispenser model?
  \item Does current magisterial judgment permit constructive use?
\end{enumerate}

Under those tests, ``maternal cooperation'' and ``maternal intercession'' are strong; ``Advocate'' is sound with its Christological subject and analogical scope stated; ``Mediatrix'' demands careful qualification and ``Mediatrix of All Graces'' should not be used as an unqualified causal claim; ``Co-redemptrix'' is always inappropriate under MPF 22. Obedience to this judgment is not diminished Marian devotion. It is devotion purified by the mystery the dogmas themselves confess: all that Mary is magnifies what God alone has done in Christ.
